%%%%%%%%%%%%%%%%
%%% lev 27 %%%
%%%%%%%%%%%%%%%%

\Chapter{27}

\Verse{1}
{
	\hP{וַיְדַבֵּר יְהֹוָה אֶל מֹשֶׁה לֵּאמֹר׃}
}
{
	\eP{And YHWH spoke to Moses, saying,}
}
{
}

\Verse{2}
{
	\hP{דַּבֵּר אֶל בְּנֵי יִשְׂרָאֵל וְאָמַרְתָּ אֲלֵהֶם אִישׁ כִּי יַפְלִא נֶדֶר בְּעֶרְכְּךָ נְפָשֹׁת לַיהֹוָה׃}
}
{
	\eP{
		“Speak to the children of Israel. And you shall say to
		them: When a man will express a vow by your appraisal of
		persons for YHWH,
	}
}
{
}

\Verse{3}
{
	\hP{וְהָיָה עֶרְכְּךָ הַזָּכָר מִבֶּן עֶשְׂרִים שָׁנָה וְעַד בֶּן שִׁשִּׁים שָׁנָה וְהָיָה עֶרְכְּךָ חֲמִשִּׁים שֶׁקֶל כֶּסֶף בְּשֶׁקֶל הַקֹּדֶשׁ׃}
}
{
	\eP{
		then your appraisal shall be: “A male from twenty years
		old up to sixty years old: then your appraisal shall be
		fifty shekels of silver by the shekel of the Holy.
	}
}
{
}

\Verse{4}
{
	\hP{וְאִם נְקֵבָה הִוא וְהָיָה עֶרְכְּךָ שְׁלֹשִׁים שָׁקֶל׃}
}
{
	\eP{
		And if it is a female, then your appraisal shall be thirty
		shekels.
	}
}
{
}

\Verse{5}
{
	\hP{וְאִם מִבֶּן חָמֵשׁ שָׁנִים וְעַד בֶּן עֶשְׂרִים שָׁנָה וְהָיָה עֶרְכְּךָ הַזָּכָר עֶשְׂרִים שְׁקָלִים וְלַנְּקֵבָה עֲשֶׂרֶת שְׁקָלִים׃}
}
{
	\eP{
		And if one is from five years old up to twenty years old,
		then your appraisal for a male shall be twenty shekels and
		for a female ten shekels.
	}
}
{
}

\Verse{6}
{
	\hP{וְאִם מִבֶּן חֹדֶשׁ וְעַד בֶּן חָמֵשׁ שָׁנִים וְהָיָה עֶרְכְּךָ הַזָּכָר חֲמִשָּׁה שְׁקָלִים כָּסֶף וְלַנְּקֵבָה עֶרְכְּךָ שְׁלֹשֶׁת שְׁקָלִים כָּסֶף׃}
}
{
	\eP{
		And if one is from a month old up to five years old, then
		your appraisal for a male shall be five shekels of silver
		and your appraisal for a female three shekels.
	}
}
{
}

\Verse{7}
{
	\hP{וְאִם מִבֶּן שִׁשִּׁים שָׁנָה וָמַעְלָה אִם זָכָר וְהָיָה עֶרְכְּךָ חֲמִשָּׁה עָשָׂר שָׁקֶל וְלַנְּקֵבָה עֲשָׂרָה שְׁקָלִים׃}
}
{
	\eP{
		And if one is from sixty years old and up, if a male then
		your appraisal shall be fifteen shekels, and for a female
		ten shekels.
	}
}
{
}

\Verse{8}
{
	\hP{וְאִם מָךְ הוּא מֵעֶרְכֶּךָ וְהֶעֱמִידוֹ לִפְנֵי הַכֹּהֵן וְהֶעֱרִיךְ אֹתוֹ הַכֹּהֵן עַל פִּי אֲשֶׁר תַּשִּׂיג יַד הַנֹּדֵר יַעֲרִיכֶנּוּ הַכֹּהֵן׃}
}
{
	\eP{
		And if he is lower than your estimate, then one shall
		stand him in front of the priest, and the priest shall
		appraise him. The priest shall appraise him on the basis
		of what the hand of the one making the vow can attain.
	}
}
{
}

\Verse{9}
{
	\hP{וְאִם בְּהֵמָה אֲשֶׁר יַקְרִיבוּ מִמֶּנָּה קרְבָּן לַיהֹוָה כֹּל אֲשֶׁר יִתֵּן מִמֶּנּוּ לַיהֹוָה יִהְיֶה קֹּדֶשׁ׃}
}
{
	\eP{
		“And if it is an animal from which people would bring an
		offering to YHWH, everything that he will give of it to
		YHWH will be a holy thing.
	}
}
{
%	\footnote{
%		an animal from which people would bring an offering. Meaning: an animal
%		that is permissible for sacrifice, as opposed to an impure animal (see
%		v. 11).
%	}
}

\Verse{10}
{
	\hP{לֹא יַחֲלִיפֶנּוּ וְלֹא יָמִיר אֹתוֹ טוֹב בְּרָע אוֹ רַע בְּטוֹב וְאִם הָמֵר יָמִיר בְּהֵמָה בִּבְהֵמָה וְהָיָה הוּא וּתְמוּרָתוֹ יִהְיֶה קֹּדֶשׁ׃}
}
{
	\eP{
		And he shall not exchange it and shall not trade it, good
		for bad or bad for good; and if he does trade an animal
		for an animal, then both it and the one traded for it
		shall be a holy thing.
	}
}
{
}

\Verse{11}
{
	\hP{וְאִם כּל בְּהֵמָה טְמֵאָה אֲשֶׁר לֹא יַקְרִיבוּ מִמֶּנָּה קרְבָּן לַיהֹוָה וְהֶעֱמִיד אֶת הַבְּהֵמָה לִפְנֵי הַכֹּהֵן׃}
}
{
	\eP{
		And if it is any impure animal, from which people would
		not bring an offering to YHWH, then he shall stand the
		animal in front of the priest,
	}
}
{
}

\Verse{12}
{
	\hP{וְהֶעֱרִיךְ הַכֹּהֵן אֹתָהּ בֵּין טוֹב וּבֵין רָע כְּעֶרְכְּךָ הַכֹּהֵן כֵּן יִהְיֶה׃}
}
{
	\eP{
		and the priest shall appraise it. Whether good or bad:
		according to your appraisal by the priest, so it shall be.
	}
}
{
}

\Verse{13}
{
	\hP{וְאִם גָּאֹל יִגְאָלֶנָּה וְיָסַף חֲמִישִׁתוֹ עַל עֶרְכֶּךָ׃}
}
{
	\eP{
		And if he will redeem it, then he shall add a fifth of it
		to your appraisal.
	}
}
{
}

\Verse{14}
{
	\hP{וְאִישׁ כִּי יַקְדִּשׁ אֶת בֵּיתוֹ קֹדֶשׁ לַיהֹוָה וְהֶעֱרִיכוֹ הַכֹּהֵן בֵּין טוֹב וּבֵין רָע כַּאֲשֶׁר יַעֲרִיךְ אֹתוֹ הַכֹּהֵן כֵּן יָקוּם׃}
}
{
	\eP{
		“And when a man will consecrate his house to be a holy
		thing to YHWH, then the priest shall appraise it. Whether
		good or bad: as the priest will appraise it, so it shall
		stand.
	}
}
{
}

\Verse{15}
{
	\hP{וְאִם הַמַּקְדִּישׁ יִגְאַל אֶת בֵּיתוֹ וְיָסַף חֲמִישִׁית כֶּסֶף עֶרְכְּךָ עָלָיו וְהָיָה לוֹ׃}
}
{
	\eP{
		And if the one who consecrates it will redeem his house,
		then he shall add a fifth of the money of your appraisal
		to it, and it shall be his.
	}
}
{
}

\Verse{16}
{
	\hP{וְאִם מִשְּׂדֵה אֲחֻזָּתוֹ יַקְדִּישׁ אִישׁ לַיהֹוָה וְהָיָה עֶרְכְּךָ לְפִי זַרְעוֹ זֶרַע חֹמֶר שְׂעֹרִים בַּחֲמִשִּׁים שֶׁקֶל כָּסֶף׃}
}
{
	\eP{
		“And if a man will consecrate any of a field of his
		possession to YHWH, then your appraisal shall be on the
		basis of its seed: the seed of a homer of barley at fifty
		shekels of silver.
	}
}
{
%	\footnote{
%		on the basis of its seed. The value of the land is determined by how
%		much seed one can plant in it.
%	}
}

\Verse{17}
{
	\hP{אִם מִשְּׁנַת הַיֹּבֵל יַקְדִּישׁ שָׂדֵהוּ כְּעֶרְכְּךָ יָקוּם׃}
}
{
	\eP{
		If he will consecrate his field from the jubilee year, it
		shall stand according to your appraisal;
	}
}
{
}

\Verse{18}
{
	\hP{וְאִם אַחַר הַיֹּבֵל יַקְדִּישׁ שָׂדֵהוּ וְחִשַּׁב לוֹ הַכֹּהֵן אֶת הַכֶּסֶף עַל פִּי הַשָּׁנִים הַנּוֹתָרֹת עַד שְׁנַת הַיֹּבֵל וְנִגְרַע מֵעֶרְכֶּךָ׃}
}
{
	\eP{
		and if he will consecrate his field after the jubilee,
		then the priest shall figure the money for him on the
		basis of the years that are left until the jubilee year,
		and it shall be subtracted from your appraisal.
	}
}
{
}

\Verse{19}
{
	\hP{וְאִם גָּאֹל יִגְאַל אֶת הַשָּׂדֶה הַמַּקְדִּישׁ אֹתוֹ וְיָסַף חֲמִשִׁית כֶּסֶף עֶרְכְּךָ עָלָיו וְקָם לוֹ׃}
}
{
	\eP{
		And if the one who consecrates it will redeem the field,
		then he shall add a fifth of the money of your appraisal
		to it, and it will stand as his.
	}
}
{
}

\Verse{20}
{
	\hP{וְאִם לֹא יִגְאַל אֶת הַשָּׂדֶה וְאִם מָכַר אֶת הַשָּׂדֶה לְאִישׁ אַחֵר לֹא יִגָּאֵל עוֹד׃}
}
{
	\eP{
		And if he will not redeem the field, or if he has sold the
		field to another man, it shall not be redeemed anymore.
	}
}
{
}

\Verse{21}
{
	\hP{וְהָיָה הַשָּׂדֶה בְּצֵאתוֹ בַיֹּבֵל קֹדֶשׁ לַיהֹוָה כִּשְׂדֵה הַחֵרֶם לַכֹּהֵן תִּהְיֶה אֲחֻזָּתוֹ׃}
}
{
	\eP{
		And when it goes out in the jubilee, the field shall be a
		holy thing to YHWH, like a devoted field; his possession
		shall be the priest’s.
	}
}
{
%	\footnote{
%		like a devoted field. Anything that one formally devotes to God cannot
%		ever be sold or redeemed, and it becomes the priests’ property (Lev
%		27:28–29; Num 18:14).
%	}
}

\Verse{22}
{
	\hP{וְאִם אֶת שְׂדֵה מִקְנָתוֹ אֲשֶׁר לֹא מִשְּׂדֵה אֲחֻזָּתוֹ יַקְדִּישׁ לַיהֹוָה׃}
}
{
	\eP{
		“And if he will consecrate any of a field that he has
		bought, that is not from a field of his possession, to
		YHWH,
	}
}
{
}

\Verse{23}
{
	\hP{וְחִשַּׁב לוֹ הַכֹּהֵן אֵת מִכְסַת הָעֶרְכְּךָ עַד שְׁנַת הַיֹּבֵל וְנָתַן אֶת הָעֶרְכְּךָ בַּיּוֹם הַהוּא קֹדֶשׁ לַיהֹוָה׃}
}
{
	\eP{
		then the priest shall figure for him the count of your
		appraisal up to the jubilee year, and he shall give your
		appraisal on that day, a holy thing to YHWH.
	}
}
{
}

\Verse{24}
{
	\hP{בִּשְׁנַת הַיּוֹבֵל יָשׁוּב הַשָּׂדֶה לַאֲשֶׁר קָנָהוּ מֵאִתּוֹ לַאֲשֶׁר לוֹ אֲחֻזַּת הָאָרֶץ׃}
}
{
	\eP{
		In the jubilee year the field shall go back to the one
		from whom it was bought, to the one to whom the possession
		of the land belongs.
	}
}
{
}

\Verse{25}
{
	\hP{וְכל עֶרְכְּךָ יִהְיֶה בְּשֶׁקֶל הַקֹּדֶשׁ עֶשְׂרִים גֵּרָה יִהְיֶה הַשָּׁקֶל׃}
}
{
	\eP{
		And all of your appraisal shall be by the shekel of the
		Holy. (The shekel shall be twenty gerah.)
	}
}
{
}

\Verse{26}
{
	\hP{אַךְ בְּכוֹר אֲשֶׁר יְבֻכַּר לַיהֹוָה בִּבְהֵמָה לֹא יַקְדִּישׁ אִישׁ אֹתוֹ אִם שׁוֹר אִם שֶׂה לַיהֹוָה הוּא׃}
}
{
	\eP{
		“Except: a firstling of the animals—which as a firstling
		is committed to YHWH—no man shall consecrate it. Whether
		an ox or a sheep, it is YHWH’s.
	}
}
{
}

\Verse{27}
{
	\hP{וְאִם בַּבְּהֵמָה הַטְּמֵאָה וּפָדָה בְעֶרְכֶּךָ וְיָסַף חֲמִשִׁתוֹ עָלָיו וְאִם לֹא יִגָּאֵל וְנִמְכַּר בְּעֶרְכֶּךָ׃}
}
{
	\eP{
		And if it is of the impure animals, then he shall redeem
		it at your appraisal, and he shall add to it a fifth of
		it. And if it will not be redeemed then it shall be sold
		at your appraisal.
	}
}
{
}

\Verse{28}
{
	\hP{אַךְ כּל חֵרֶם אֲשֶׁר יַחֲרִם אִישׁ לַיהֹוָה מִכּל אֲשֶׁר לוֹ מֵאָדָם וּבְהֵמָה וּמִשְּׂדֵה אֲחֻזָּתוֹ לֹא יִמָּכֵר וְלֹא יִגָּאֵל כּל חֵרֶם קֹדֶשׁ קדָשִׁים הוּא לַיהֹוָה׃}
}
{
	\eP{
		“Except: any devoted thing that a man will devote to YHWH
		from anything that he has—from human or animal or from a
		field of his possession—shall not be sold and shall not be
		redeemed. Any devoted thing: it is a holy of holies to
		YHWH.
	}
}
{
}

\Verse{29}
{
	\hP{כּל חֵרֶם אֲשֶׁר יחֳרַם מִן הָאָדָם לֹא יִפָּדֶה מוֹת יוּמָת׃}
}
{
	\eP{
		Anyone who will be devoted from humans shall not be
		redeemed. He shall be put to death.
	}
}
{
}

\Verse{30}
{
	\hP{וְכל מַעְשַׂר הָאָרֶץ מִזֶּרַע הָאָרֶץ מִפְּרִי הָעֵץ לַיהֹוָה הוּא קֹדֶשׁ לַיהֹוָה׃}
}
{
	\eP{
		“And all the tithe of the land, from the land’s seed, from
		the tree’s fruit: it is YHWH’s, a holy thing to YHWH.
	}
}
{
}

\Verse{31}
{
	\hP{וְאִם גָּאֹל יִגְאַל אִישׁ מִמַּעַשְׂרוֹ חֲמִשִׁיתוֹ יֹסֵף עָלָיו׃}
}
{
	\eP{
		And if a man will redeem some of his tithe, he shall add
		to it a fifth of it.
	}
}
{
}

\Verse{32}
{
	\hP{וְכל מַעְשַׂר בָּקָר וָצֹאן כֹּל אֲשֶׁר יַעֲבֹר תַּחַת הַשָּׁבֶט הָעֲשִׂירִי יִהְיֶה קֹּדֶשׁ לַיהֹוָה׃}
}
{
	\eP{
		And all the tithe of a herd or a flock, everything that
		passes under the rod: the tenth shall be a holy thing to
		YHWH.
	}
}
{
}

\Verse{33}
{
	\hP{לֹא יְבַקֵּר בֵּין טוֹב לָרַע וְלֹא יְמִירֶנּוּ וְאִם הָמֵר יְמִירֶנּוּ וְהָיָה הוּא וּתְמוּרָתוֹ יִהְיֶה קֹּדֶשׁ לֹא יִגָּאֵל׃}
}
{
	\eP{
		He shall not inspect it either for good or for bad. And he
		shall not trade it. And if he will trade it, then both it
		and the one traded for it shall be a holy thing. It shall
		not be redeemed.”
	}
}
{
}

\Verse{34}
{
	\hP{אֵלֶּה הַמִּצְוֺת אֲשֶׁר צִוָּה יְהֹוָה אֶת מֹשֶׁה אֶל בְּנֵי יִשְׂרָאֵל בְּהַר סִינָי׃}
}
{
	\eP{
		These are the commandments that YHWH commanded Moses for
		the children of Israel in Mount Sinai.
	}
}
{
%	\footnote{
%		These are the commandments that YHWH commanded Moses for the children of
%		Israel in Mount Sinai. Literary studies of Leviticus are likely to begin
%		by raising the question of whether this book that contains so much law
%		and so little narrative belongs in a literary study at all. But at
%		bottom we do not have to argue that Leviticus is great literature in
%		order to justify its inclusion in a literary analysis. It is not
%		literature, not in the way in which almost anyone ever means that term.
%		But the literary study of the Hebrew Bible has to come to terms with
%		Leviticus, for it is embedded in the biblical narrative. Indeed, the
%		fact that no law code from ancient Israel survived independently, that
%		law codes survived only in contexts of narrative, is notable in itself.
%		Leviticus’s laws are integral—and assumed thereafter. The book begins
%		with the notation that YHWH speaks these things to Moses from the
%		Tabernacle (1:1), and it ends here with the reminder that “These are the
%		commandments that YHWH commanded Moses for the children of Israel at
%		Mount Sinai.” Some of the laws presume the preceding narrative, e.g.,
%		“You shall not do like what is done in the land of Egypt, in which you
%		lived …” (18:3), which presumes the bondage in Egypt. Legal authority is
%		vested in the priests (the power to teach and enforce, 10:10–11), which
%		is an essential element in narratives in subsequent books. As new
%		political figures, especially kings, will rise, the compatibility or
%		conflict between them and the priests will be critical to the narrative,
%		as it was critical in history. Leviticus thus is a design for an
%		organized society of people who help one another, who do not
%		intentionally injure one another, who respect one another’s property and
%		relationships, who regularly assemble to celebrate together, who
%		acknowledge their errors and atone for them, who regard life—in humans
%		and in animals—as sacred, who pursue purity in various forms, who
%		respect law, and who are utterly loyal to one God. It is against the
%		background of this design that the book of Numbers must be understood.
%		COMMENTARY ON THE TORAH
%		The book of Numbers is the story of a journey.
%		Against the backdrop of the book of Leviticus, with its absence of
%		movement, Numbers is entirely about movement. The journey as a literary
%		theme has been a recurring component of world literature from the epic
%		of Gilgamesh, the oldest known book, and the Odyssey to Alice’s
%		Adventures in Wonderland and The Wonderful Wizard of Oz. It reflects
%		real experience, it is a ready metaphor for human lives, and it is a
%		connector between the world of literature and the world of dreams, in
%		which experiences of journeys are common. The journey in literature can
%		be the experience of an individual or a group. In the book of Numbers it
%		is both. It is the people of Israel’s journey from Sinai to the border
%		of Canaan, from the site of the establishment of the covenant to the
%		place where it is to be fulfilled. And it is Moses’ journey from Mount
%		Sinai to Mount Abarim, from the place where he has seen God to the place
%		of his death. One must read the story in Numbers with a sense of what is
%		pictured: a great mass of people (603,550 adult males, plus the junior
%		and elderly males and the females of all ages), an entire nation, in
%		movement—with a visible miracle (the column of cloud and fire) and
%		miraculous feeding (manna) happening in their presence at all times over
%		a span of nearly forty years. It is the only age in the Bible in which
%		miracle is a daily fact of life. There is an organization for travel, a
%		hierarchy for community life (Moses, judges, priests), and formal
%		rituals of time (Sabbath, holidays, and new moons) and of space
%		(sacrifice, incense, vestments, the Tabernacle, fringes on clothing).
%		There is a destination anticipated in a future in which some of them
%		will participate and some will not. They are in an interim condition,
%		conscious of the life in Egypt behind them and the promised life in
%		front of them. Their mentality is a mentality of slaves. There are
%		various groups and individuals among them with an interest in positions
%		of leadership. On the evidence of the episodes of the story, they are to
%		be pictured as deeply afraid. Against this setting, the deity is
%		pictured frequently as speaking and acting in response to human events
%		rather than by divine initiation. Divine communications are fewer and
%		briefer than in the preceding books. Leviticus is primarily concerned
%		with God’s giving the rules to the people. Now Numbers is concerned with
%		their first experiences living under these rules and with God’s
%		responses to the courses of action they take. Against the backdrop of
%		the commandments of Leviticus, Numbers is the story of a people coming
%		to terms with having a constitution of laws, and coming to terms with
%		their relationship of holiness with the deity. And in the middle of this
%		stands Moses. The story is about the people and God, but it is focused
%		more than ever on Moses, his personality and especially his conflicts.
%	}
}
